\section{Розділ~\ref{sec:glutcomputer}. Що обчислюють \nt{GLUT}-нейрони}

Логічні твердження розділу --- теореми про схеми з невід'ємних ваг, виведені в самому розділі; дослід підтверджує модель нейрона й те, що схеми, зібрані за цими теоремами (детектор збігів, лінія затримки, самопідтримувана група, ланцюжок), у мозку справді трапляються.

{\footnotesize
\setlength{\tabcolsep}{3pt}\renewcommand{\arraystretch}{1.15}
\begin{longtable}{>{\raggedright}p{0.5cm}>{\raggedright}p{4.0cm}>{\raggedright}p{5.6cm}>{\raggedright}p{2.3cm}>{\raggedright\arraybackslash}p{1.7cm}}
\toprule
№ & Твердження & Дослід і що виміряно & Джерело & Статус \\
\midrule\endfirsthead
\toprule
№ & Твердження & Дослід і що виміряно & Джерело & Статус \\
\midrule\endhead
1 & Нейрон --- RC-коло з порогом і скиданням~\eqref{eq:glutneuron} & У тіло пірамідних нейронів кори щура (зрізи) подавали шумовий струм, схожий на синаптичний потік у живому мозку. Залежність середньої частоти від середнього й розкиду струму описується інтегрувальним нейроном із витоком, якщо додати повільну адаптацію частоти. Діапазони $\tau$ і затримок наведено в розділі без посилання на дослід & \cite{Rauch2003} & виміряно \\
2 & Модель~\eqref{eq:glutneuron} --- спрощення: гілки входів самі видають місцеві імпульси & Запис безпосередньо з дендритів пірамідних нейронів зорової кори миші in vivo: зоровий подразник викликає місцеві дендритні імпульси, що залежать від рецепторів NMDA; їх придушення розширює налаштування нейрона на орієнтацію & \cite{Smith2013} & виміряно \\
3 & Вікно збігу $\Delta_{\max}=\tau\ln\frac{w}{\theta-w}$~\eqref{eq:window}; при $\theta=1.5w$ воно дорівнює $0.7\tau$ & Наслідок моделі~\eqref{eq:glutneuron}. Пряме вимірювання вузького вікна підсумовування є для нейронів зі швидким гальмуванням (розділ~\ref{sec:gaba}, рядок~3 його таблиці) & висновок у розділі & теорія \\
4 & Мережа з самих лише \nt{GLUT}-нейронів обчислює лише монотонні функції входів (твердження~\ref{prop:monotone}) & Доведення в розділі: ітерація від тиші з неспадною правою частиною. Це твердження про модель; досліду, який перевіряв би його на живій мережі без гальмування, немає & висновок у розділі & теорія \\
5 & Органи чуття дають пару проводів: одні клітини відповідають на посилення подразника, інші --- на послаблення & Сітківка: уже на біполярних клітинах сигнал колбочок розходиться на вмикальний і вимикальний канали. Фармакологічний блок вмикальних біполярних клітин у мавпи: канали йдуть окремо до кори; після блоку погіршується виявлення посилення світла, але не послаблення & \cite{Schiller1992} & виміряно \\
6 & З двопровідним кодом мережа з \nt{GLUT}-нейронів обчислює будь-яку логічну функцію асинхронно, без спільного такту, ціною вдвічі більшої кількості нейронів (твердження~\ref{prop:dualrail}, \eqref{eq:dualrail}, рис.~\ref{fig:dualxor}) & Двопровідний код і визначення готовності за самим сигналом --- стандартна техніка асинхронних схем, нечутливих до затримок. Схему виключного АБО з шести нейронів зібрано в розділі. Досліду, який показував би, що мозок обчислює логічні функції саме двопровідним кодом, немає & \cite{Sparso2001} & теорія \\
7 & Найбільша різниця часу приходу звуку у вуха в людини $\approx0.6\ms$, а мозок розрізняє близько десяти мікросекунд & Слухачі в досліді з двома варіантами відповіді: поріг розрізнення різниці часу (75\,\% правильних) --- $9$\,мкс для шуму в смузі $150$--$1700$\,Гц, $11$\,мкс для тону $1$\,кГц, $28$\,мкс для клацання. Межа $0.6\ms$ --- розрахунок розділу, $0.22\,\text{м}/343\,\text{м/с}$ & \cite{KlumppEady1956} & виміряно \\
8 & У птахів різниця часів перетворюється на місце через лінії затримки й детектори збігів~\eqref{eq:itd} (рис.~\ref{fig:itd}) & Сипуха: аксони від двох вух входять у ядро детекторів збігів із протилежних боків; час приходу імпульсів уздовж аксона змінюється впорядковано з глибиною, і розкид затримок крізь ядро ($160$\,мкс на $700$\,мкм) збігається з діапазоном міжвушних затримок & \cite{Carr1990} & виміряно \\
9 & У ссавців ряду затримок не знайшли; ту саму задачу розв'язують детектори збігів із точно розрахованим гальмуванням від \nt{GLYC}-нейронів & Запис in vivo в медіальній верхній оливі піщанки: відповіді не утворюють карти напрямків; підведення гліцину та його блокатора через мікропіпетку показує, що точно розраховане в часі \emph{гліцинове} гальмування задає робочий діапазон затримок. Гальмування тут від \nt{GLYC}, а не від \nt{GABA} & \cite{Brand2002,Grothe2010} & виміряно \\
10 & Група із сильним зворотним зв'язком --- тригер; самопідтримувана активність триває секунди після подразника (рис.~\ref{fig:bistable}) & Префронтальна кора мавпи в задачі з відкладеним рухом очей до запам'ятованої точки (затримка $1$--$6$\,с): $87$ із $170$ нейронів, пов'язаних із задачею, змінюють частоту під час затримки, у $79\,\%$ із них це залежить від напрямку на запам'ятовану точку. Що цю активність тримає зворотний зв'язок із двома стійкими станами, --- теорія & \cite{Funahashi1989,Wang2001} & виміряно \\
11 & Біт записується, якщо приплив триває довше $\approx0.7\tau$ (рис.~\ref{fig:recurrentgroup}б) & Розрахунок рівняння $\tau\,\dd r/\dd t=-r+f(wr+I)$ методом Ейлера при $w=1$, $I=0.6$; скрипту в \texttt{models/} немає. Досліду немає & розрахунок у розділі & модель \\
12 & Пам'ять можна зберігати в полегшенні синапсів, майже без імпульсів & Теорія: залишковий кальцій у закінченнях тримає запис близько секунди. Опора: у медіальній префронтальній корі тхора (запис із кількох нейронів одночасно) є підмережа пірамідних нейронів, зв'язаних полегшуваними синапсами з тривалою післядією. Огляд спостережень «тихих» проміжків під час утримання в пам'яті. Який спосіб кора використовує коли --- відкрите питання & \cite{Mongillo2008,WangMarkram2006,Stokes2015} & гіпотеза \\
13 & Пам'ять стирається втомою: запас медіатора виснажується & Парні записи пірамідних нейронів кори: відповідь синапсу падає за частих імпульсів, швидкість падіння задається ймовірністю викиду. Що саме це гасить самопідтримувану групу, безпосередньо не показано & \cite{Tsodyks1997} & виміряно \\
14 & Ланцюжок груп --- таймер, що запускається подією; у співочих птахів нейрони спрацьовують суворо по черзі & Запис упізнаних нейронів вищого вокального центру зебрової амадини уві сні й під час співу: кожен нейрон, що посилає вихід у моторне ядро, видає одну коротку пачку в один точний момент пісні, і нейрони спрацьовують послідовно & \cite{Hahnloser2002} & виміряно \\
\bottomrule
\end{longtable}}
